\section{Introduction}
\label{sec:introduction}

Petrophysical property estimation is a central task in reservoir
characterization. Wireline logs are dense along depth and can be acquired in
most wells, but properties such as porosity, clay volume, or shale volume are
usually calibrated by core or laboratory measurements that are sparse and
expensive \citep{ellis2007well,asquith2004basic}. This mismatch creates a
recurring inverse problem: a model must infer a continuous target-property
profile from dense indirect logs while honoring a small number of direct
measurements in the target well.

Machine-learning regressors are attractive in this setting because they can
learn nonlinear relations between logs and target properties from previously
interpreted wells \citep{bishop2006pattern}. However, a direct map from logs to
targets may underuse sparse core observations available in the target well. A
common direct workaround is to concatenate the sparse measurement vector with
the logs and train a supervised model from \([u,b]\) to \(y\). Although this can
be competitive, it treats the sparse measurements as features rather than as
measurement evidence about the reconstructed signal.

Compressed sensing (CS) offers the complementary perspective: measurements
should constrain the reconstruction through a forward model
\citep{donoho2006compressed,candes2006robust,candes2006stable}. Classical CS
usually assumes sparsity in a fixed basis, such as wavelets or discrete cosine
atoms. In cross-well petrophysical prediction, this assumption can be too
restrictive. Geological variability is not necessarily sparse in a universal
basis, and fixed-basis residual correction may fail to outperform strong
supervised baselines.

Compressed sensing with generative models (CSGM) replaces fixed-basis sparsity
with a learned generative range \citep{bora2017compressed}. Instead of seeking
a sparse coefficient vector, one searches for a latent code \(z\) such that the
generated signal \(G(z)\) matches the measurements. This is well aligned with
petrophysical windows: the decoder can represent plausible target profiles,
while sparse core observations can select the profile that is consistent with
the target-well evidence.

This paper proposes Conditional Latent-Prior CSGM (CLP-CSGM) for
petrophysical property estimation from dense logs and sparse core measurements.
The method trains an autoencoder decoder \(G\) on target-property windows and a
ridge prior \(h\) that maps the dense log window \(u\) to an initial latent
code \(z_0=h(u)\). At test time, the final latent code is obtained by solving a
measurement-consistent inverse problem,
\begin{equation}
  \hat z =
  \arg\min_z
  \|M G(z)-b\|_2^2
  + \lambda \|z-z_0(u)\|_2^2,
  \qquad
  \hat y = G(\hat z).
  \label{eq:intro_clp_csgm}
\end{equation}
The first term penalizes sparse core inconsistency; the second term keeps the
solution close to the log-conditioned latent prior.

The contributions are:
\begin{itemize}
  \item We formulate CLP-CSGM, a conditional generative compressed sensing
  method for petrophysical property windows.
  \item We provide a theoretical interpretation of the objective as a
  Tikhonov/MAP latent inverse problem and relate it to standard CSGM guarantees.
  \item We evaluate the method in low-data cross-well clay-volume prediction
  and in a real-well F03-4 GR-only porosity stress test.
  \item We show that CLP-CSGM Ridge improves over strong direct baselines in
  sparse-measurement regimes, and we use structural ablations to show that both
  the conditional latent prior and the sparse-measurement term are needed.
\end{itemize}

The remainder of the paper reviews related work, defines CLP-CSGM, describes the
datasets and experimental protocol, reports quantitative and ablation results,
and discusses the regimes in which the method is most appropriate.
